\documentclass[../../main.tex]{subfiles}
\begin{document}

{\bf [解]}

$10000=2^4\cdot5^4$である．$a^2-a=a(a-1)$において$a$と$a-1$は互いに素である．
また題意より$a$は奇数だから，$a$は$5^4$の倍数，$a-1$は$2^4$の倍数であることが必要である．
したがって$m,n\in\mathbb N$として
\begin{align}
  a=m\cdot5^4=n\cdot2^4+1 \label{eq:1}
\end{align}
とおける．$5^4=625=39\cdot16+1$より，
\begin{align}
  m(39\cdot16+1)-n\cdot16=1 \quad \therefore 16(39m-n)=1-m \label{eq:2}
\end{align}
となる．よって$1-m$は$16$の倍数であり，$k\in\mathbb Z$として
\begin{align}
  m=1-16k \label{eq:3}
\end{align}
である．これを\cref{eq:1}に代入して$3\le a\le9999$から，
\begin{align}
  3\le 625(1-16k)\le 9999
\end{align}
を得る．これを満たす整数$k$は$k=0$だけだから，$m=1$であり，よって\cref{eq:1}から
\begin{align}
  a=625
\end{align}
である．逆に，$625$は奇数で，
\begin{align}
  625^2-625=625\cdot624=10000\cdot 39
\end{align}
だから，確かに条件を満たして十分である．

以上から求めるのは
\begin{align}
  a = 625
\end{align}
である．

\end{document}
